%!TeX root=../houndtop.tex
\chapter{The Stapletons of Merripit House}
\lettrine[lines=4]{T}{he} fresh beauty of the following morning did something to efface from our minds the grim and gray impression which had been left upon both of us by our first experience of Baskerville Hall. As Sir Henry and I sat at breakfast the sunlight flooded in through the high mullioned windows, throwing watery patches of colour from the coats of arms which covered them. The dark panelling glowed like bronze in the golden rays, and it was hard to realize that this was indeed the chamber which had struck such a gloom into our souls upon the evening before.

»I guess it is ourselves and not the house that we have to blame!« said the baronet. »We were tired with our journey and chilled by our drive, so we took a gray view of the place. Now we are fresh and well, so it is all cheerful once more.«

»And yet it was not entirely a question of imagination,« I an\-swered. »Did you, for example, happen to hear someone, a woman I think, sobbing in the night?«

»That is curious, for I did when I was half asleep fancy that I heard something of the sort. I waited quite a time, but there was no more of it, so I concluded that it was all a dream.«

»I heard it distinctly, and I am sure that it was really the sob of a woman.«

»We must ask about this right away.« He rang the bell and asked Barrymore whether he could account for our experience. It seemed to me that the pallid features of the butler turned a shade paler still as he listened to his master's question.

»There are only two women in the house, Sir Henry,« he an\-swered. »One is the scullery-maid, who sleeps in the other wing. The other is my wife, and I can answer for it that the sound could not have come from her.«

And yet he lied as he said it, for it chanced that after breakfast I met Mrs Barrymore in the long corridor with the sun full upon her face. She was a large, impassive, heavy-featured woman with a stern set expression of mouth. But her tell-tale eyes were red and glanced at me from between swollen lids. It was she, then, who wept in the night, and if she did so her husband must know it. Yet he had taken the obvious risk of discovery in declaring that it was not so. Why had he done this? And why did she weep so bitterly? Already round this pale-faced, handsome, black-bearded man there was gathering an atmosphere of mystery and of gloom. It was he who had been the first to discover the body of Sir Charles, and we had only his word for all the circumstances which led up to the old man's death. Was it possible that it was Barrymore after all whom we had seen in the cab in Regent Street? The beard might well have been the same. The cabman had described a somewhat shorter man, but such an impression might easily have been erroneous. How could I settle the point forever? Obviously the first thing to do was to see the Grimpen postmaster, and find whether the test telegram had really been placed in Barrymore's own hands. Be the answer what it might, I should at least have something to report to Sherlock Holmes.

Sir Henry had numerous papers to examine after breakfast, so that the time was propitious for my excursion. It was a pleasant walk of four miles along the edge of the moor, leading me at last to a small gray hamlet, in which two larger buildings, which proved to be the inn and the house of Dr Mortimer, stood high above the rest. The postmaster, who was also the village grocer, had a clear recollection of the telegram.

»Certainly, sir,« said he, »I had the telegram delivered to Mr Barrymore exactly as directed.«

»Who delivered it?«

»My boy here. James, you delivered that telegram to Mr Barrymore at the Hall last week, did you not?«

»Yes, father, I delivered it.«

»Into his own hands?« I asked.

»Well, he was up in the loft at the time, so that I could not put it into his own hands, but I gave it into Mrs Barrymore's hands, and she promised to deliver it at once.«

»Did you see Mr Barrymore?«

»No, sir; I tell you he was in the loft.«

»If you didn't see him, how do you know he was in the loft?«

»Well, surely his own wife ought to know where he is,« said the postmaster testily. »Didn't he get the telegram? If there is any mistake it is for Mr Barrymore himself to complain.«

It seemed hopeless to pursue the inquiry any farther, but it was clear that in spite of Holmes's ruse we had no proof that Barrymore had not been in London all the time. Suppose that it were so\allowbreak—\allowbreak suppose that the same man had been the last who had seen Sir Charles alive, and the first to dog the new heir when he returned to England. What then? Was he the agent of others or had he some sinister design of his own? What interest could he have in persecuting the Baskerville family? I thought of the strange warning clipped out of the leading article of the \textit{Times}. Was that his work or was it possibly the doing of someone who was bent upon counteracting his schemes? The only conceivable motive was that which had been suggested by Sir Henry, that if the family could be scared away a comfortable and permanent home would be secured for the Barrymores. But surely such an explanation as that would be quite inadequate to account for the deep and subtle scheming which seemed to be weaving an invisible net round the young baronet. Holmes himself had said that no more complex case had come to him in all the long series of his sensational investigations. I prayed, as I walked back along the gray, lonely road, that my friend might soon be freed from his preoccupations and able to come down to take this heavy burden of responsibility from my shoulders.

Suddenly my thoughts were interrupted by the sound of running feet behind me and by a voice which called me by name. I turned, expecting to see Dr Mortimer, but to my surprise it was a stranger who was pursuing me. He was a small, slim, clean-shaven, prim-faced man, flaxen-haired and lean-jawed, between thirty and forty years of age, dressed in a gray suit and wearing a straw hat. A tin box for botanical specimens hung over his shoulder and he carried a green butterfly-net in one of his hands.

\begin{figure}[tbhp]
\centering
\includegraphics[width=\linewidth]{07_stranger}
\caption{It was a stranger pursuing me.}
\end{figure}

»You will, I am sure, excuse my presumption, Dr Watson,« said he, as he came panting up to where I stood. »Here on the moor we are homely folk and do not wait for formal introductions. You may possibly have heard my name from our mutual friend, Mortimer. I am Stapleton, of Merripit House.«

»Your net and box would have told me as much,« said I, »for I knew that Mr Stapleton was a naturalist. But how did you know me?«

»I have been calling on Mortimer, and he pointed you out to me from the window of his surgery as you passed. As our road lay the same way I thought that I would overtake you and introduce myself. I trust that Sir Henry is none the worse for his journey?«

»He is very well, thank you.«

»We were all rather afraid that after the sad death of Sir Charles the new baronet might refuse to live here. It is asking much of a wealthy man to come down and bury himself in a place of this kind, but I need not tell you that it means a very great deal to the country-side. Sir Henry has, I suppose, no superstitious fears in the matter?«

»I do not think that it is likely.«

»Of course you know the legend of the fiend dog which haunts the family?«

»I have heard it.«

»It is extraordinary how credulous the peasants are about here! Any number of them are ready to swear that they have seen such a creature upon the moor.« He spoke with a smile, but I seemed to read in his eyes that he took the matter more seriously. »The story took a great hold upon the imagination of Sir Charles, and I have no doubt that it led to his tragic end.«

»But how?«

»His nerves were so worked up that the appearance of any dog might have had a fatal effect upon his diseased heart. I fancy that he really did see something of the kind upon that last night in the Yew Alley. I feared that some disaster might occur, for I was very fond of the old man, and I knew that his heart was weak.«

»How did you know that?«

»My friend Mortimer told me.«

»You think, then, that some dog pursued Sir Charles, and that he died of fright in consequence?«

»Have you any better explanation?«

»I have not come to any conclusion.«

»Has Mr Sherlock Holmes?«

The words took away my breath for an instant, but a glance at the placid face and steadfast eyes of my companion showed that no surprise was intended.

»It is useless for us to pretend that we do not know you, Dr Watson,« said he. »The records of your detective have reached us here, and you could not celebrate him without being known yourself. When Mortimer told me your name he could not deny your identity. If you are here, then it follows that Mr Sherlock Holmes is interesting himself in the matter, and I am naturally curious to know what view he may take.«

»I am afraid that I cannot answer that question.«

»May I ask if he is going to honour us with a visit himself?«

»He cannot leave town at present. He has other cases which engage his attention.«

»What a pity! He might throw some light on that which is so dark to us. But as to your own researches, if there is any possible way in which I can be of service to you I trust that you will command me. If I had any indication of the nature of your suspicions or how you propose to investigate the case, I might perhaps even now give you some aid or advice.«

»I assure you that I am simply here upon a visit to my friend, Sir Henry, and that I need no help of any kind.«

»Excellent!« said Stapleton. »You are perfectly right to be wary and discreet. I am justly reproved for what I feel was an unjustifiable intrusion, and I promise you that I will not mention the matter again.«

We had come to a point where a narrow grassy path struck off from the road and wound away across the moor. A steep, boulder-sprinkled hill lay upon the right which had in bygone days been cut into a granite quarry. The face which was turned towards us formed a dark cliff, with ferns and brambles growing in its niches. From over a distant rise there floated a gray plume of smoke.

»A moderate walk along this moor-path brings us to Merripit House,« said he. »Perhaps you will spare an hour that I may have the pleasure of introducing you to my sister.«

My first thought was that I should be by Sir Henry's side. But then I remembered the pile of papers and bills with which his study table was littered. It was certain that I could not help with those. And Holmes had expressly said that I should study the neighbours upon the moor. I accepted Stapleton's invitation, and we turned together down the path.

»It is a wonderful place, the moor,« said he, looking round over the undulating downs, long green rollers, with crests of jagged granite foaming up into fantastic surges. »You never tire of the moor. You cannot think the wonderful secrets which it contains. It is so vast, and so barren, and so mysterious.«

»You know it well, then?«

»I have only been here two years. The residents would call me a newcomer. We came shortly after Sir Charles settled. But my tastes led me to explore every part of the country round, and I should think that there are few men who know it better than I do.«

»Is it hard to know?«

»Very hard. You see, for example, this great plain to the north here with the queer hills breaking out of it. Do you observe anything remarkable about that?«

»It would be a rare place for a gallop.«

»You would naturally think so and the thought has cost several their lives before now. You notice those bright green spots scattered thickly over it?«

»Yes, they seem more fertile than the rest.«

Stapleton laughed.


\begin{figure}[tbhp]
\centering
\includegraphics[width=\linewidth]{07_greatgrimpenmire}
\caption{»That is the great Grimpen Mire,« said he.}
\end{figure}

»That is the great Grimpen Mire,« said he. »A false step yonder means death to man or beast. Only yesterday I saw one of the moor ponies wander into it. He never came out. I saw his head for quite a long time craning out of the bog-hole, but it sucked him down at last. Even in dry seasons it is a danger to cross it, but after these autumn rains it is an awful place. And yet I can find my way to the very heart of it and return alive. By George, there is another of those miserable ponies!«

Something brown was rolling and tossing among the green sedges. Then a long, agonized, writhing neck shot upward and a dreadful cry echoed over the moor. It turned me cold with horror, but my companion's nerves seemed to be stronger than mine.

»It's gone!« said he. »The mire has him. Two in two days, and many more, perhaps, for they get in the way of going there in the dry weather, and never know the difference until the mire has them in its clutches. It's a bad place, the great Grimpen Mire.«

»And you say you can penetrate it?«

»Yes, there are one or two paths which a very active man can take. I have found them out.«


%\afterpage{\clearpage}

»But why should you wish to go into so horrible a place?«

»Well, you see the hills beyond? They are really islands cut off on all sides by the impassable mire, which has crawled round them in the course of years. That is where the rare plants and the butterflies are, if you have the wit to reach them.«

»I shall try my luck some day.«

He looked at me with a surprised face.

»For God's sake put such an idea out of your mind,« said he. »Your blood would be upon my head. I assure you that there would not be the least chance of your coming back alive. It is only by remembering certain complex landmarks that I am able to do it.«

»Halloa!« I cried. »What is that?«

A long, low moan, indescribably sad, swept over the moor. It filled the whole air, and yet it was impossible to say whence it came. From a dull murmur it swelled into a deep roar, and then sank back into a melancholy, throbbing murmur once again. Stapleton looked at me with a curious expression in his face.

»Queer place, the moor!« said he.

»But what is it?«

»The peasants say it is the Hound of the Baskervilles calling for its prey. I've heard it once or twice before, but never quite so loud.«

I looked round, with a chill of fear in my heart, at the huge swelling plain, mottled with the green patches of rushes. Nothing stirred over the vast expanse save a pair of ravens, which croaked loudly from a tor behind us.

»You are an educated man. You don't believe such nonsense as that?« said I. »What do you think is the cause of so strange a sound?«

»Bogs make queer noises sometimes. It's the mud settling, or the water rising, or something.«

»No, no, that was a living voice.«

»Well, perhaps it was. Did you ever hear a bittern booming?«

»No, I never did.«

%»It's a very rare bird—practically extinct—in England now, but all things are possible upon the moor. Yes, I should not be surprised to learn that what we have heard is the cry of the last of the bitterns.\footnote{The Eurasian bittern, \textit{Botaurus stellaris}, is still threatened in the UK but is not extinct. In 2022 the Royal Society for the Protection of Birds announced that the British bittern population had, for the first time, surpassed its estimated 1880s population.}«

»It's a very rare bird—practically extinct—in England now, but all things are possible upon the moor. Yes, I should not be surprised to learn that what we have heard is the cry of the last of the bitterns.«


»It's the weirdest, strangest thing that ever I heard in my life.«

»Yes, it's rather an uncanny place altogether. Look at the hill- side yonder. What do you make of those?«

The whole steep slope was covered with gray circular rings of stone, a score of them at least.

»What are they? Sheep-pens?«

»No, they are the homes of our worthy ancestors. Prehistoric man lived thickly on the moor, and as no one in particular has lived there since, we find all his little arrangements exactly as he left them. These are his wigwams with the roofs off. You can even see his hearth and his couch if you have the curiosity to go inside.«

»But it is quite a town. When was it inhabited?«

%»Neolithic man—no date.\footnote{Modern estimate for the hut circles on Dartmoor is circa second millennium BCE.}«

»Neolithic man—no date.«

»What did he do?«

%»He grazed his cattle on these slopes, and he learned to dig for tin when the bronze sword began to supersede the stone axe. Look at the great trench in the opposite hill. That is his mark. Yes, you will find some very singular points about the moor, Dr Watson. Oh, excuse me an instant! It is surely Cyclopides\footnote{Genus of butterflies of the family commonly called »skippers« for their erratic flight.}.«

»He grazed his cattle on these slopes, and he learned to dig for tin when the bronze sword began to supersede the stone axe. Look at the great trench in the opposite hill. That is his mark. Yes, you will find some very singular points about the moor, Dr Watson. Oh, excuse me an instant! It is surely Cyclopides.«


A small fly or moth had fluttered across our path, and in an instant Stapleton was rushing with extraordinary energy and speed in pursuit of it. To my dismay the creature flew straight for the great mire, and my acquaintance never paused for an instant, bounding from tuft to tuft behind it, his green net waving in the air. His gray clothes and jerky, zigzag, irregular progress made him not unlike some huge moth himself. I was standing watching his pursuit with a mixture of admiration for his extraordinary activity and fear lest he should lose his footing in the treacherous mire, when I heard the sound of steps, and turning round found a woman near me upon the path. She had come from the direction in which the plume of smoke indicated the position of Merripit House, but the dip of the moor had hid her until she was quite close.

I could not doubt that this was the Miss Stapleton of whom I had been told, since ladies of any sort must be few upon the moor, and I remembered that I had heard someone describe her as being a beauty. The woman who approached me was certainly that, and of a most uncommon type. There could not have been a greater contrast between brother and sister, for Stapleton was neutral tinted, with light hair and gray eyes, while she was darker than any brunette whom I have seen in England—slim, elegant, and tall. She had a proud, finely cut face, so regular that it might have seemed impassive were it not for the sensitive mouth and the beautiful dark, eager eyes. With her perfect figure and elegant dress she was, indeed, a strange apparition upon a lonely moorland path. Her eyes were on her brother as I turned, and then she quickened her pace towards me. I had raised my hat and was about to make some explanatory remark, when her own words turned all my thoughts into a new channel.

»Go back!« she said. »Go straight back to London, instantly.«

\makeatletter
\@ifclasswith{scrbook}{a5paper}
{%
	\begin{figure}[tbhp]
	\centering
	\includegraphics[width=\linewidth]{07_goback}
	\caption{»Go back!« she said.}
	\end{figure}
}{%

}
\makeatother



I could only stare at her in stupid surprise. Her eyes blazed at me, and she tapped the ground impatiently with her foot.

»Why should I go back?« I asked.

»I cannot explain.« She spoke in a low, eager voice, with a curious lisp in her utterance. »But for God's sake do what I ask you. Go back and never set foot upon the moor again.«

»But I have only just come.«

»Man, man!« she cried. »Can you not tell when a warning is for your own good? Go back to London! Start to-night! Get away from this place at all costs! Hush, my brother is coming! Not a word of what I have said. Would you mind getting that orchid for me among the mares-tails yonder? We are very rich in orchids on the moor, though, of course, you are rather late to see the beauties of the place.«

Stapleton had abandoned the chase and came back to us breathing hard and flushed with his exertions.

»Halloa, Beryl!« said he, and it seemed to me that the tone of his greeting was not altogether a cordial one.

»Well, Jack, you are very hot.«

»Yes, I was chasing a Cyclopides. He is very rare and seldom found in the late autumn. What a pity that I should have missed him!« He spoke unconcernedly, but his small light eyes glanced incessantly from the girl to me.

»You have introduced yourselves, I can see.«

»Yes. I was telling Sir Henry that it was rather late for him to see the true beauties of the moor.«

»Why, who do you think this is?«

»I imagine that it must be Sir Henry Baskerville.«

»No, no,« said I. »Only a humble commoner, but his friend. My name is Dr Watson.«

A flush of vexation passed over her expressive face. »We have been talking at cross purposes,« said she.

»Why, you had not very much time for talk,« her brother remarked with the same questioning eyes.

»I talked as if Dr Watson were a resident instead of being merely a visitor,« said she. »It cannot much matter to him whether it is early or late for the orchids. But you will come on, will you not, and see Merripit House?«

A short walk brought us to it, a bleak moorland house, once the farm of some grazier in the old prosperous days, but now put into repair and turned into a modern dwelling. An orchard surrounded it, but the trees, as is usual upon the moor, were stunted and nipped, and the effect of the whole place was mean and melancholy. We were admitted by a strange, wizened, rusty-coated old manservant, who seemed in keeping with the house. Inside, however, there were large rooms furnished with an elegance in which I seemed to recognize the taste of the lady. As I looked from their windows at the interminable granite-flecked moor rolling unbroken to the farthest horizon I could not but marvel at what could have brought this highly educated man and this beautiful woman to live in such a place.

»Queer spot to choose, is it not?« said he as if in answer to my thought. »And yet we manage to make ourselves fairly happy, do we not, Beryl?«

»Quite happy,« said she, but there was no ring of conviction in her words.

»I had a school,« said Stapleton. »It was in the north country. The work to a man of my temperament was mechanical and uninteresting, but the privilege of living with youth, of helping to mould those young minds, and of impressing them with one's own character and ideals, was very dear to me. However, the fates were against us. A serious epidemic broke out in the school and three of the boys died. It never recovered from the blow, and much of my capital was irretrievably swallowed up. And yet, if it were not for the loss of the charming companionship of the boys, I could rejoice over my own misfortune, for, with my strong tastes for botany and zoology, I find an unlimited field of work here, and my sister is as devoted to Nature as I am. All this, Dr Watson, has been brought upon your head by your expression as you surveyed the moor out of our window.«

»It certainly did cross my mind that it might be a little dull—less for you, perhaps, than for your sister.«

»No, no, I am never dull,« said she, quickly.

»We have books, we have our studies, and we have interesting neighbours. Dr Mortimer is a most learned man in his own line. Poor Sir Charles was also an admirable companion. We knew him well, and miss him more than I can tell. Do you think that I should intrude if I were to call this afternoon and make the acquaintance of Sir Henry?«

»I am sure that he would be delighted.«

%»Then perhaps you would mention that I propose to do so. We may in our humble way do something to make things more easy for him until he becomes accustomed to his new surroundings. Will you come upstairs, Dr Watson, and inspect my collection of Lepidoptera\footnote{The order of insects including moths and butterflies.}? I think it is the most complete one in the south-west of England. By the time that you have looked through them lunch will be almost ready.«

»Then perhaps you would mention that I propose to do so. We may in our humble way do something to make things more easy for him until he becomes accustomed to his new surroundings. Will you come upstairs, Dr Watson, and inspect my collection of Lepidoptera? I think it is the most complete one in the south-west of England. By the time that you have looked through them lunch will be almost ready.«

But I was eager to get back to my charge. The melancholy of the moor, the death of the unfortunate pony, the weird sound which had been associated with the grim legend of the Baskervilles, all these things tinged my thoughts with sadness. Then on the top of these more or less vague impressions there had come the definite and distinct warning of Miss Stapleton, delivered with such intense earnestness that I could not doubt that some grave and deep reason lay behind it. I resisted all pressure to stay for lunch, and I set off at once upon my return journey, taking the grass-grown path by which we had come.

It seems, however, that there must have been some short cut for those who knew it, for before I had reached the road I was astounded to see Miss Stapleton sitting upon a rock by the side of the track. Her face was beautifully flushed with her exertions, and she held her hand to her side.

»I have run all the way in order to cut you off, Dr Watson,« said she. »I had not even time to put on my hat. I must not stop, or my brother may miss me. I wanted to say to you how sorry I am about the stupid mistake I made in thinking that you were Sir Henry. Please forget the words I said, which have no application whatever to you.«

»But I can't forget them, Miss Stapleton,« said I. »I am Sir Hen\-ry's friend, and his welfare is a very close concern of mine. Tell me why it was that you were so eager that Sir Henry should return to London.«

»A woman's whim, Dr Watson. When you know me better you will understand that I cannot always give reasons for what I say or do.«

»No, no. I remember the thrill in your voice. I remember the look in your eyes. Please, please, be frank with me, Miss Stapleton, for ever since I have been here I have been conscious of shadows all round me. Life has become like that great Grimpen Mire, with little green patches everywhere into which one may sink and with no guide to point the track. Tell me then what it was that you meant, and I will promise to convey your warning to Sir Henry.«

An expression of irresolution passed for an instant over her face, but her eyes had hardened again when she answered me.

»You make too much of it, Dr Watson,« said she. »My brother and I were very much shocked by the death of Sir Charles. We knew him very intimately, for his favourite walk was over the moor to our house. He was deeply impressed with the curse which hung over the family, and when this tragedy came I naturally felt that there must be some grounds for the fears which he had expressed. I was distressed therefore when another member of the family came down to live here, and I felt that he should be warned of the danger which he will run. That was all which I intended to convey.«

»But what is the danger?«

»You know the story of the hound?«

\begin{figure}[tbhp]
\centering
\includegraphics[width=\linewidth]{07_storyhound}
\caption{»You know the story of the hound?«}
\end{figure}

»I do not believe in such nonsense.«

»But I do. If you have any influence with Sir Henry, take him away from a place which has always been fatal to his family. The world is wide. Why should he wish to live at the place of danger?«

»Because it \textit{is} the place of danger. That is Sir Henry's nature. I fear that unless you can give me some more definite information than this it would be impossible to get him to move.«

»I cannot say anything definite, for I do not know anything definite.«

»I would ask you one more question, Miss Stapleton. If you meant no more than this when you first spoke to me, why should you not wish your brother to overhear what you said? There is nothing to which he, or anyone else, could object.«

»My brother is very anxious to have the Hall inhabited, for he thinks it is for the good of the poor folk upon the moor. He would be very angry if he knew that I have said anything which might induce Sir Henry to go away. But I have done my duty now and I will say no more. I must get back, or he will miss me and suspect that I have seen you. Good-bye!« She turned and had disappeared in a few minutes among the scattered boulders, while I, with my soul full of vague fears, pursued my way to Baskerville Hall.

%\cleardoubleevenpage
%\vfill
%\begin{figure}[tbph]
%\centering
%\includegraphics[width=\linewidth]{thespot}
%\caption{He took us both to show us the spot}
%\end{figure}
%\vfill
%\thispagestyle{empty}
%\clearpage
